%%
%%Huidig onboardingsproces
%%

\chapter{Huidige onboarding- en offboardingprocessen}
\label{ch:HuidigeProcessen}

In dit hoofdstuk wordt de huidige werkwijze besproken die gebruikt wordt binnen het gemeentebestuur van Berlare.

\section{Beschrijving van de huidige processen}

De beschrijving van de huidige processen binnen de gemeente Berlare.

\subsection{Huidig onboardingsproces}

Het huidige onboardingsproces binnen het gemeentebestuur van Berlare omvat verschillende administratieve en technische stappen om de toegang van een nieuwe medewerker tot interne systemen op te bouwen.

\subsubsection{Algemene werking}

Het bestaande onboardingproces van nieuwe medewerkers binnen het gemeentebestuur werd geanalyseerd om de huidige werkwijze, betrokken diensten en uitgevoerde stappen in kaart te brengen.

Het proces omvat administratieve, technische en communicatieve handelingen die voornamelijk door de personeelsdienst en de ICT-dienst worden uitgevoerd. De personeelsdienst verzamelt de noodzakelijke persoonsgegevens en functiegegevens en bezorgt deze aan de ICT-dienst, die vervolgens de benodigde accounts en toegangen configureert.

Figuur~\ref{fig:HuidigeOnboarding} geeft een vereenvoudigd overzicht van de huidige onboardingflow binnen het gemeentebestuur.

\begin{figure}[!h]
    \centering
    \includegraphics[width=0.75\textwidth]{HuidigeOnboarding.png}
    \caption[Huidig onboardingsproces]{Overzicht van het huidige onboardingproces en de interactie tussen de betrokken diensten}
    \label{fig:HuidigeOnboarding}
\end{figure}

\subsubsection{Huidige werkwijze binnen de ICT-dienst}

Binnen de ICT-dienst verloopt een groot deel van het onboardingproces momenteel manueel. Gebruikersaccounts worden individueel aangemaakt binnen Active Directory, waarna gebruikers handmatig toegevoegd worden aan relevante beveiligingsgroepen en distributielijsten.

Daarnaast worden mailboxen geconfigureerd en worden rechten ingesteld voor verschillende interne systemen en gedeelde mappen. Ook licenties worden afzonderlijk toegekend afhankelijk van de functie en de benodigde softwaretoegang van de medewerker.

Ook de communicatie met de nieuwe medewerker gebeurt grotendeels manueel. Welkomstmails met handleidingen en documentatie worden afzonderlijk verzonden. Daarnaast wordt een tijdelijke wachtwoordbrief gegenereerd, afgedrukt en fysiek bezorgd.

\subsubsection{Fysieke configuratiehandelingen}

Naast digitale configuraties omvat het onboardingproces eveneens verschillende fysieke handelingen binnen de ICT-dienst. Hierbij gaat het onder andere om het voorbereiden van hardware, het installeren van toestellen en het configureren van softwarepakketten.

Deze handelingen worden momenteel individueel uitgevoerd afhankelijk van de functie en de noden van de nieuwe medewerker. Hierdoor bestaat het onboardingproces uit een combinatie van digitale configuraties en fysieke installatietaken.

\subsection{Huidig offboardingsproces}

Het huidige offboardingsproces binnen het gemeentebestuur van Berlare omvat verschillende administratieve en technische stappen om de toegang van een vertrekkende medewerker tot interne systemen af te bouwen.

\subsubsection{Algemene werking}

De personeelsdienst of het diensthoofd meldt de uitdiensttreding aan de ICT-dienst. Vervolgens worden verschillende configuraties handmatig aangepast, waaronder het uitschakelen van gebruikersaccounts, verwijderen van toegangsrechten en aanpassen van gedeelde postvakken.

Omdat deze handelingen afhankelijk zijn van de einddatum van het contract, worden ze doorgaans op of kort na de laatste werkdag uitgevoerd

Figuur \ref{fig:HuidigeOffboarding} toont een visuele voorstelling van het huidige offboardingsproces.

\begin{figure}[!h]
    \centering
    \includegraphics[width=0.75\textwidth]{HuidigeOffboarding.png}
    \caption[Huidig offboardingsproces]{Overzicht van het huidige offboardingproces}
    \label{fig:HuidigeOffboarding}
\end{figure}

\subsubsection{Werkwijze binnen de ICT-dienst}

Binnen de ICT-dienst verloopt het offboardingproces momenteel grotendeels manueel. Medewerkers moeten handmatig gebruikersaccounts opzoeken binnen Active Directory en deze vervolgens uitschakelen of aanpassen. Daarnaast worden gebruikers verwijderd uit beveiligingsgroepen, distributielijsten en Microsoft 365-groepen. Ook mailboxrechten op gedeelde postvakken moeten afzonderlijk verwijderd worden.

De mailbox van een voormalige medewerker wordt omgezet naar een gedeeld postvak zodat historische communicatie behouden blijft. Daarnaast wordt een Out of Office ingesteld met informatie over de opvolging van de mailbox. Ook de toegekende Microsoft 365-licenties worden verwijderd zodat deze opnieuw beschikbaar zijn.

Verder moeten van de laptops de gebruikersaccounts verwijderd worden. Ook op de servers moeten deze opnieuw verwijderd worden. Op deze manier verdwijnen de bestanden die lokaal op laptops opgeslagen waren, waardoor er dus ruimte vrij komt op een laptop en deze dus opnieuw gebruikt kan worden. 

\subsubsection{Beheer van toegangen en mailboxen}

Een belangrijk onderdeel van het offboardingproces bestaat uit het correct verwijderen van toegangsrechten. Wanneer accounts of rechten niet tijdig verwijderd worden, bestaat het risico dat voormalige medewerkers nog toegang behouden tot interne systemen of gevoelige gegevens.

Omdat toegangsrechten verspreid beheerd worden over verschillende platformen, waaronder Active Directory, Entra ID en Exchange Online, vereist dit proces momenteel meerdere afzonderlijke handelingen binnen verschillende beheertools.

Daarnaast bestaat er momenteel geen centrale opvolging waarmee eenvoudig gecontroleerd kan worden welke onderdelen van het offboardingproces reeds uitgevoerd werden.

\section{Communicatie en samenwerking tussen diensten}

Voor de hudige processen binnen het gemeentebestuur van Berlare is er ook samenwerking en communicatie nodig tussen de interne diensten.

\subsection{Afhankelijkheden tussen diensten}

Binnen het huidige onboardingproces zijn meerdere diensten afhankelijk van elkaar om bepaalde stappen correct te kunnen uitvoeren. De personeelsdienst levert de noodzakelijke gegevens aan voor het aanmaken of uitschakelen van gebruikersaccounts, terwijl andere interne diensten afhankelijk zijn van het beschikbaar zijn van een e-mailadres of gebruikersaccount om toegang te voorzien tot dienstspecifieke toepassingen.

Hierdoor ontstaat een proces waarbij verschillende stappen pas uitgevoerd kunnen worden nadat voorgaande configuraties succesvol voltooid zijn. Een correcte communicatie tussen diensten speelt daarom een belangrijke rol binnen het onboardingproces.

\subsection{Terugkoppeling van uitgevoerde taken}

Binnen de huidige processen gebeurt de terugkoppeling van uitgevoerde taken voornamelijk manueel. Medewerkers communiceren onderling wanneer bepaalde configuraties of administratieve stappen voltooid zijn.

Andere diensten zijn afhankelijk van meldingen of terugkoppelingen wanneer gebruikersaccounts succesvol aangemaakt werden. Deze informatie is noodzakelijk om bijkomende configuraties uit te voeren binnen dienstspecifieke softwaretoepassingen.

Deze werkwijze maakt het moeilijker om de voortgang van het onboardingproces centraal op te volgen en kan leiden tot onduidelijkheid over de status van bepaalde stappen.

\section{Knelpunten binnen de huidige processen}

Binnen het gemeentebestuur van Berlare zijn er een aantal duidelijke knelpunten binnen beide processen.

\subsection{Manuele configuraties}

Een groot deel van het onboarding- en offboardingproces wordt momenteel manueel uitgevoerd. Het aanmaken en uitschakelen van gebruikersaccounts, het beheren van groepslidmaatschappen en licenties en het configureren van mailboxen vereisen verschillende repetitieve handelingen van de ICT-dienst.

Daarnaast verhoogt manuele configuratie de kans op inconsistenties en menselijke fouten. Foutieve groepslidmaatschappen, ontbrekende rechten of vergeten configuraties kunnen leiden tot bijkomende correcties en vertragingen. Verder kan het niet verwijderen van gebruikers doordat dit vergeten wordt, leiden tot veiligheidsrisico's.

\subsection{Gebrek aan centrale opvolging}

Een bijkomend knelpunt is het ontbreken van een centrale opvolging van uitgevoerde taken. Omdat verschillende stappen verspreid uitgevoerd worden door meerdere medewerkers en diensten, ontstaat er soms onduidelijkheid over de voortgang van het proces. Hierdoor ontstaan afhankelijkheden waarbij medewerkers moeten navragen of bepaalde configuraties reeds voltooid zijn.

\subsection{Complexiteit van onboarding en offboarding}

Beide processen bestaan uit verschillende onderling verbonden stappen waarbij meerdere systemen betrokken zijn. Naast gebruikersbeheer in Active Directory moeten eveneens mailboxen, licenties, distributielijsten en interne toepassingen correct geconfigureerd of aangepast worden.

Daarnaast verhoogt de hybride infrastructuur, waarbij zowel Active Directory als Microsoft 365 en Entra ID gebruikt worden, de complexiteit van beide processen. Synchronisaties tussen lokale en cloudomgevingen zorgen ervoor dat bepaalde configuraties pas uitgevoerd kunnen worden nadat eerdere stappen succesvol voltooid werden.

Deze combinatie van afhankelijkheden, verschillende beheersystemen en manuele configuraties bemoeilijkt een efficiënte opvolging van onboarding en offboarding binnen de huidige werking van het gemeentebestuur.
